\section{RNN 在计算机视觉中的应用}

\subsection{图像描述生成（Image Captioning）}

最经典的 CNN + RNN 应用是\textbf{图像描述生成}：用 CNN 提取图像特征，用 RNN 生成自然语言描述。

\begin{figure}[H]
\centering
\includegraphics[width=0.95\textwidth]{slides-images/slide-037.jpg}
\caption{CNN + RNN 图像描述模型：CNN（如 ImageNet 预训练的 ResNet）提取图像特征作为 RNN 的初始输入，RNN 逐步生成描述文字\protect\footnotemark}
\end{figure}
\footnotetext{Slides 第38页。}

\begin{figure}[H]
\centering
\includegraphics[width=0.95\textwidth]{slides-images/slide-038.jpg}
\caption{更具体的 CNN+RNN 架构：使用 CNN 倒数第二层的特征向量作为 RNN 隐藏状态的附加输入\protect\footnotemark}
\end{figure}
\footnotetext{Slides 第39页。}

\begin{figure}[H]
\centering
\includegraphics[width=0.95\textwidth]{slides-images/slide-040.jpg}
\caption{图像描述的成功案例和失败案例\protect\footnotemark}
\end{figure}
\footnotetext{Slides 第41页。}

\begin{warningbox}{共现偏置与幻觉}
早期的图像描述模型存在严重的\textbf{共现偏置}问题：
\begin{itemize}
    \item 看到海滩就假设有冲浪板
    \item 看到手持物体就猜测是常见搭配
    \item 看到棒球运动就默认是``投球''而非``接球''
\end{itemize}
这种\textbf{幻觉}（hallucination）问题在今天的视觉语言模型中依然普遍存在。其根源在于模型学习了训练数据中的统计关联，而非真正理解场景内容。
\end{warningbox}

\subsection{视觉问答（Visual Question Answering）}

\begin{figure}[H]
\centering
\includegraphics[width=0.95\textwidth]{slides-images/slide-042.jpg}
\caption{视觉问答（VQA）：给定图像和问题，模型输出答案。可以通过序列生成或多选分类来实现\protect\footnotemark}
\end{figure}
\footnotetext{Slides 第43页。}

\subsection{其他应用}

RNN 在计算机视觉中的其他应用包括：
\begin{itemize}
    \item \textbf{视觉对话}（Visual Dialogue）：围绕图像进行多轮对话
    \item \textbf{视觉导航}（Visual Navigation）：根据视觉输入生成移动指令序列
\end{itemize}

\subsection{本章小结}

RNN 与 CNN 的结合开创了视觉-语言多模态模型的先河。尽管这些方法已被 Transformer 基础的模型所取代，但 CNN+RNN 的架构思想——将视觉编码器与序列生成器组合——在今天的视觉语言模型（如 LLaVA、GPT-4V）中以不同形式延续。

%% ========== 多层 RNN ==========

